\section{Resilience 与 Public Service：事件之后如何继续工作}

安全事故和法律程序会影响员工、研究人员、用户与领导者的健康、身份和职业。\term{resilience}（韧性）不是假装没有伤害，而是通过支持、恢复、反思和重新参与，逐步重建可持续服务能力。课程谈到 Sullivan 在乌克兰公益项目中的工作；这部分不能作为案件证据，却提供了一个重要的人类系统视角。

\subsection{Human recovery 也是 incident lifecycle}

传统 postmortem 关注 root cause、control gap 和 owner，却可能忽略长期值班、公开压力、道德伤害和家庭负担。组织如果只要求“恢复服务”，会让关键人员在下一次事件前耗尽。Recovery plan 应包括轮换、心理支持、法律与 HR 协助、休假、角色调整以及重新进入工作的路径。

\lecturefigure{14-resilience-loop.png}{韧性来自支持、恢复、意义重建与重新服务的循环。}{本地字幕 27:00--29:20；概念重绘。}

\begin{knowledgebox}{读图：恢复不是回到事故前的状态}
Incident/case 造成业务与个人冲击；support 提供 family、peer、professional 和 legal care；recovery 需要 rest、perspective 与 new routines；service/learning 把经验转化为 community work、mentoring 和 better systems。箭头不是保证每个人都会线性恢复，而是提醒组织为不同速度保留空间。将 resilience 解释为“个人要更坚强”会掩盖制度责任。
\end{knowledgebox}

\begin{warningbox}{不要把公益或个人成长当成法律洗白}
一个人在事件后的善行可以解释其价值选择和恢复过程，却不能自动改变已经裁判的事实；同样，法律责任也不应让组织忽视健康与人性。课程讲义应允许两种判断同时存在，而不是用励志叙事取代证据分析。
\end{warningbox}

\subsection{Government technical talent}

公共部门需要理解 cloud、identity、AI、supply chain 和 incident operations 的技术人才，否则规则可能只描述目标而无法评估实现。技术领导者参与 government 不只意味着全职任职，也包括 standards work、advisory groups、公开评论、incident coordination 和短期服务。参与者必须处理利益冲突，并把公共责任置于公司 lobbying 之外。

\lecturefigure{15-government-talent.png}{技术人才只有留在政策、实施与行业反馈闭环中，才能提高公共治理质量。}{本地字幕 36:20--39:17；概念重绘。}

\begin{knowledgebox}{读图：人才流动要形成制度记忆}
Technical expertise 帮助描述 systems、threats 和 operational cost；policy design 把目标写成 rules、incentives 与 oversight；implementation 建立 agency capability、metrics 和 enforcement；industry feedback 返回 evidence、failure modes 与 burden。若专家只短期“空降”而没有文档、career path 和 permanent staff，知识会随个人离开；若行业只提供立场，不提供可验证数据，也无法形成高质量规则。
\end{knowledgebox}

\teachervoice{讲者最后呼吁技术人员不要把 government 当成别人的问题。课堂提示：当软件系统承载公共生活时，懂系统的人需要帮助设计公共约束；否则规则仍会出现，只是更可能在信息不足和危机压力下出现。}

\subsection{本章小结}

Resilience 把个人恢复纳入长期安全能力，public service 把技术经验送回制度反馈。两者都要求从英雄叙事转向支持系统和组织记忆。

